\section{原则五：Prompt Engineering 工具描述}

这是最后一个、也是\textbf{最有效}的方法：对工具的描述和规格（descriptions and specs）进行 Prompt Engineering。因为这些描述会被加载到 Agent 的上下文中，它们可以系统性地引导 Agent 走向有效的工具调用行为。

\subsection{编写工具描述的核心原则}

编写工具描述时，应该像\textbf{向团队新人介绍你的工具}一样思考。考虑你可能隐含带入的上下文——专用查询格式、小众术语的定义、底层资源之间的关系——并将这些\textbf{显式化}。

具体建议：

\begin{enumerate}[leftmargin=1.5em]
    \item \textbf{消除歧义}：用严格的数据模型明确描述（并强制执行）预期的输入和输出
    \item \textbf{参数命名清晰}：不要用 \texttt{user} 这样的模糊名称，使用 \texttt{user\_id} 明确指代
    \item \textbf{提供充足的上下文}：专用格式、术语定义、资源关系都应在描述中说明
    \item \textbf{结合评估验证}：即使是微小的描述优化也能带来显著的性能提升
\end{enumerate}

\begin{importantbox}{实战案例：SWE-bench 上的突破}
Claude Sonnet 3.5 在 SWE-bench Verified 评估上取得了当时的 state-of-the-art 性能，而关键手段之一就是\textbf{对工具描述进行精确的优化}。这些优化大幅降低了错误率并提高了任务完成率，证明了 prompt engineering 工具描述的巨大潜力。
\end{importantbox}

\subsection{进阶资源}

Anthropic 还推荐了以下进阶内容：

\begin{itemize}[leftmargin=1.5em]
    \item \textbf{Developer Guide}：更多工具定义的最佳实践
    \item \textbf{工具动态加载机制}：了解工具如何被动态加载到 Claude 的 system prompt 中
    \item \textbf{Tool Annotations}（MCP Server）：声明哪些工具需要开放世界访问或执行破坏性操作
\end{itemize}

\subsection{本章小结}

工具描述的 Prompt Engineering 是最具杠杆效应的优化手段。像对待新人入职文档一样编写工具描述，将隐式知识显式化，消除歧义，并通过评估验证每次优化的效果。即使微小的改进也可能带来戏剧性的性能提升。


